\documentclass[11pt]{article}

\usepackage[T1]{fontenc}
\usepackage[utf8]{inputenc}
\usepackage[english]{babel}
\usepackage[margin=1in]{geometry}
\usepackage{lmodern}
\usepackage{microtype}
\usepackage{mathtools}
\usepackage{amsmath,amssymb,amsthm}
\usepackage{array}
\usepackage{booktabs}
\usepackage{tabularx}
\usepackage{longtable}
\usepackage{float}
\usepackage{graphicx}
\usepackage{tikz}
\usetikzlibrary{arrows.meta,positioning,calc}
\usepackage{pdflscape}
\usepackage{caption}
\usepackage{enumitem}
\usepackage{fancyhdr}
\usepackage[numbers]{natbib}
\usepackage{doi}
\usepackage{xurl}
\usepackage{xr-hyper}
\usepackage{hyperref}
\externaldocument[][nocite]{build/supplement}[supplement.pdf]
\usepackage{cleveref}
\usepackage{orcidlink}
\usepackage{titling}
\usepackage{aliascnt}

% Shared manuscript macros.

%--- cluster names -------------------------------------------------------------
\newcommand{\clusterA}{Temporal Currencies}
\newcommand{\clusterB}{Moving and Self-Generated Obstructions}
\newcommand{\clusterC}{Defect Integration and Emergent Transport}
\newcommand{\clusterD}{Productive Obstructions and Sparse Saturation}

%--- series identities -----------------------------------------------------------
\newcommand{\FVI}{Foundations~VI}
\newcommand{\FoundI}{Foundations~I}
\newcommand{\FoundII}{Foundations~II}
\newcommand{\FoundIII}{Foundations~III}
\newcommand{\FoundIV}{Foundations~IV}
\newcommand{\FoundV}{Foundations~V}

%--- law names -------------------------------------------------------------------
\newcommand{\lawGOne}{Hidden Amortized Solvency}
\newcommand{\lawGTwo}{Transfinite Escrow}
\newcommand{\lawGThree}{Amortized Potential Currency}
\newcommand{\lawGFour}{Moving-Cover Exhaustion}
\newcommand{\lawGFive}{Carry-Horizon Confinement}
\newcommand{\lawGSix}{Endogenous Needle Generation}
\newcommand{\lawGSeven}{Adversarial Mobility Confinement}
\newcommand{\lawGEight}{Odometer Abelianization}
\newcommand{\lawGNine}{Defect-Evacuation-to-Transport}
\newcommand{\lawGTen}{Lossful Boundary Persistence}
\newcommand{\lawGEleven}{Local-Rule Global-Anti-Symmetry}
\newcommand{\lawGTwelve}{Finite Witness Radiation}
\newcommand{\lawGThirteen}{Thin-Orbit Saturation}

%--- one-line origin after a result ------------------------------------------------
\newcommand{\origin}[1]{\par\noindent{\small\emph{Origin.}\ #1}\par\medskip}
%--- one-line formal status after a law ----------------------------------------------
\newcommand{\formal}[1]{\par\noindent{\small\emph{In Lean.}\ #1}\par\medskip}

% Generated by scripts/paper_set_release.py.
% The date is the UTC date when this submission was assembled.
\newif\ifpaperhasdoi
\paperhasdoitrue
\newcommand{\paperdoi}{10.5281/zenodo.22254289}
\newcommand{\manuscriptversion}{Preprint v2.0}
\newcommand{\manuscriptdate}{27 September 2026}
\newcommand{\manuscriptyear}{2026}


% One shared counter, section-scoped.  Each sibling environment gets an alias counter so
% that cleveref names it by its own type; numbering is unchanged.
\newtheorem{theorem}{Theorem}[section]
\newaliascnt{lemma}{theorem}\newtheorem{lemma}[lemma]{Lemma}\aliascntresetthe{lemma}
\newaliascnt{proposition}{theorem}\newtheorem{proposition}[proposition]{Proposition}\aliascntresetthe{proposition}
\newaliascnt{corollary}{theorem}\newtheorem{corollary}[corollary]{Corollary}\aliascntresetthe{corollary}
\newaliascnt{schema}{theorem}\newtheorem{schema}[schema]{Template}\aliascntresetthe{schema}
\theoremstyle{definition}
\newaliascnt{definition}{theorem}\newtheorem{definition}[definition]{Definition}\aliascntresetthe{definition}
\newaliascnt{remark}{theorem}\newtheorem{remark}[remark]{Remark}\aliascntresetthe{remark}
\newaliascnt{example}{theorem}\newtheorem{example}[example]{Example}\aliascntresetthe{example}
\theoremstyle{plain}\newtheorem*{headline}{Theorem}

\urlstyle{same}
\captionsetup{
  font=small,
  labelfont=bf,
  labelsep=period
}
\captionsetup[table]{skip=6pt,position=top}
\crefname{suppsection}{Supplement Section}{Supplement Sections}\Crefname{suppsection}{Supplement Section}{Supplement Sections}
\crefname{theorem}{Theorem}{Theorems}
\Crefname{theorem}{Theorem}{Theorems}
\crefname{lemma}{Lemma}{Lemmas}
\Crefname{lemma}{Lemma}{Lemmas}
\crefname{proposition}{Proposition}{Propositions}
\Crefname{proposition}{Proposition}{Propositions}
\crefname{corollary}{Corollary}{Corollaries}
\Crefname{corollary}{Corollary}{Corollaries}
\crefname{schema}{Template}{Templates}
\Crefname{schema}{Template}{Templates}
\crefname{definition}{Definition}{Definitions}
\Crefname{definition}{Definition}{Definitions}
\crefname{remark}{Remark}{Remarks}
\Crefname{remark}{Remark}{Remarks}
\crefname{example}{Example}{Examples}
\Crefname{example}{Example}{Examples}
\setlist[enumerate]{leftmargin=*, itemsep=2pt, topsep=4pt}
\setlength{\droptitle}{-3.2em}
\pretitle{\begin{center}\LARGE\bfseries}
\posttitle{\par\end{center}\vspace{0.5em}}
\preauthor{\begin{center}\normalsize}
\postauthor{\par\end{center}\vspace{-0.4em}}
\predate{\begin{center}\small}
\postdate{\par\end{center}\vspace{0.6em}}
\setlength{\emergencystretch}{2em}
\clubpenalty=10000
\widowpenalty=10000
\displaywidowpenalty=10000
\interfootnotelinepenalty=10000
\allowdisplaybreaks[2]
\fancypagestyle{firstpage}{\fancyhf{}
  \fancyfoot[C]{\scriptsize Automorph Inc., Wilmington, DE, USA \quad \texttt{ioannis@automorph.io} \quad \textcopyright\ \manuscriptyear \quad CC-BY 4.0}
  \renewcommand{\headrulewidth}{0pt}
  \renewcommand{\footrulewidth}{0pt}
}

\title{Six Birds Foundations VI:\\
A Catalog of Dynamical Structural Laws}
\author{Ioannis Tsiokos\,\orcidlink{0009-0009-7659-5964}}
\date{Version 2: \manuscriptdate\\[2pt] \footnotesize (Version 1: 2 September 2026)%
  \ifpaperhasdoi\\\small\href{https://doi.org/\paperdoi}{doi:\paperdoi}\fi}
\hypersetup{
  pdftitle={Six Birds Foundations VI: A Catalog of Dynamical Structural Laws},
  pdfauthor={Ioannis Tsiokos},
  pdfsubject={Six Birds Theory, Foundations VI. \manuscriptversion, \manuscriptdate.},
  pdfkeywords={emergence calculus; dynamical structural law; liveness; template; Collatz affine ledger; split pair; reverse-and-add; Ducci map; abelian sandpile; odometer; Lean 4; Six Birds Theory},
  hidelinks
}

\begin{document}

\maketitle
\thispagestyle{firstpage}

\begin{abstract}
A static structural law classifies fixed data. Many questions about dynamics concern runs
instead: whether every orbit eventually descends, whether a run terminates while its
presentation changes, whether a cover must move with its target. This paper states thirteen
candidate laws of this kind and sorts them by what is proved. Seven are laws, each with a concrete
system in which the difficult hypothesis is proved (for one of them only a small arithmetic
example); two are general theorems with imported classical instances; four are templates, whose
implication is proved while the difficult hypothesis remains an open obligation. Lean coverage is reported separately. The concrete results below are
all checked in Lean. For the accelerated Collatz map every orbit satisfies an exact affine ledger, and at every depth $m$ two starting values agree modulo $2^m$ but
have different next division exponents, so no fixed residue class predicts the ledger. Under binary
reverse-and-add, the orbit of $22$ enters a four-phase family and never reaches a palindrome. The
Ducci map on binary vectors of length $2^k$ is nilpotent. For move systems whose distinct legal
moves commute, started where no infinite legal run exists, any two legal runs that end in stable
states end in the same state with the same firing counts. Classical results are imported with
attribution, nothing is labelled new, and the Collatz, Erd\H{o}s--Straus, Lychrel and Recam\'an
problems remain open.

\end{abstract}

\medskip
\noindent\textbf{Keywords.} emergence calculus; dynamical structural law; liveness; template;
Collatz affine ledger; split pair; reverse-and-add; Ducci map; abelian sandpile; odometer;
Six Birds Theory.

\section{Introduction}\label{sec:intro}

A structural law, in the sense this series gives the term, is a theorem about declared abstract
data --- a carrier, a quotient, a ledger, a family of patches --- so that any system supplying the
data inherits the conclusion. \FoundIV{} collected fifty-two such laws
\citep{Tsiokos2026FoundIV}. Every one of them is about fixed data: one map, one pair of quotients,
one declared finite family of patches. None quantifies over the steps of a run. Yet many of the
questions that make dynamics hard are questions about runs. Does every orbit of a map eventually
descend? Does a run terminate when the way its states are written down changes at every step?
Can a family of targets be covered when no fixed finite package of patches suffices? What happens
when the obstruction to progress is produced by the run's own history, or chosen by an adversary?
When do route choices integrate into a conserved ledger, when do local defects turn into motion,
what survives repeated loss, and how do local rules force global structure?

This paper states thirteen candidate laws, G1--G13, one for each of these questions, and sorts
them by what has actually been proved. Each entry has a general statement whose \emph{hard
hypothesis} --- a coverage property, a closed family of states, a hierarchy of forced scales ---
carries the difficult content. Seven entries are \emph{laws}: for each, the hard hypothesis is
proved here for a concrete system, although for G13 that system is only a small arithmetic example.
Two are \emph{general theorems}, whose hypotheses are ordinary data and whose classical instances
are imported. Four are \emph{templates}: the implication is proved, but its hard hypothesis is not
established in this paper for any system of interest and remains an open obligation. \Cref{tab:catalog} records the sorting, and
separately what is checked in Lean.

\subsection{A worked example: residues do not predict the Collatz ledger}\label{sec:intro-example}

The accelerated Collatz map sends an odd number $n$ to $T(n) = (3n+1)/2^{a(n)}$, where $a(n)$ is
the exponent of the largest power of $2$ dividing $3n+1$. Take $n = 5$ and $n = 9$. Both are
congruent to $1$ modulo $4$, so they agree in their last two binary digits. But $3\cdot 5 + 1 =
16$ gives $a(5) = 4$, while $3 \cdot 9 + 1 = 28$ gives $a(9) = 2$. Knowing $n$ modulo $4$ does not
determine even the first division exponent.

The same happens at every depth: for each $m \ge 1$ there are two odd numbers that agree modulo
$2^m$ whose first division exponents are $m$ and at least $m+1$ (\Cref{thm:split-pair}). The pair
above is the case $m = 2$. At the same time the bookkeeping of an orbit is exact. After $k$ steps,
with $A_k$ the sum of the first $k$ exponents,
\[
  2^{A_k}\, n_k \;=\; 3^k\, n + B_k ,
\]
for an explicit integer $B_k$ determined by the exponents (\Cref{thm:affine-ledger}). The orbit has
dropped below its start at step $k$ exactly when $2^{A_k} n > 3^k n + B_k$, an inequality between
integers. Whether every orbit eventually satisfies it is the Collatz conjecture, which this paper
does not settle. The two facts show where the difficulty lives: in an exact ledger whose next entry
is not predicted by any fixed number of low-order digits.

\subsection{Obstructions at every depth}\label{sec:intro-contribution}

Two notions recur throughout the paper. A \emph{quotient} $q_m$ records what an observer sees at
depth $m$, for example the last $m$ binary digits of a number, and a \emph{readout} $r$ records a
future fact, for example the next division exponent. A \emph{split pair} for $q_m$ and $r$ is a pair
of states with the same quotient and different readouts. If a split pair exists at every depth, no
observer who looks at a fixed number of digits can predict the readout.

The Collatz pairs above are such an obstruction for the last $m$ binary digits and the first
division exponent. The orbits of $2^L - 1$ give a second one for the time of first descent: they lie
in the residue class of $-1$ modulo $2^m$ once $L \ge m$, and they stay above their start for at
least $L - 1$ steps (\Cref{rem:collatz-paper}). So no bound on the time of first descent can depend
only on the residue of $n$ modulo a fixed power of $2$.

Each entry of the catalog is stated in one form: the data of a run, a certificate in the run's own
terms, and the named hypothesis that makes the certificate available. When that hypothesis is not
yet proved for any system of interest, the entry is a \emph{template}, and its value is that it
names the obligation precisely and lets it be recognized in unrelated systems. The split-pair
condition above, for instance, is exactly what must be avoided by any proof that a fixed-depth
reading decides a run; the same condition reappears for binary reverse-and-add, where a different
tool, a family of states closed under the map and free of palindromes, settles one orbit
(Result~B below).

\subsection{Main results}\label{sec:intro-results}

The following four results are proved in the text and checked in Lean. Their relation to the
catalog differs.
\begin{itemize}
  \item Result~A (\Cref{sec:G1}) gives the certificate that G1 asks for, in exact integer form, and
    shows that no fixed number of binary digits predicts it. It does not establish the hypotheses
    of G1.
  \item Result~B (\Cref{sec:G5}) establishes the hard hypothesis of G5 for one orbit.
  \item Result~C (\Cref{sec:G10}) is a collapse example for G10. It shows that parity cannot be the
    protected invariant for vectors of length $2^k$.
  \item Result~D (\Cref{sec:G8}) is the general statement of G8, with two examples that separate
    it from weaker properties.
\end{itemize}

\begin{headline}[A: the Collatz ledger and its split pairs]
Every orbit of the accelerated Collatz map satisfies $2^{A_k} n_k = 3^k n_0 + B_k$ for every
$k$. For every depth $m \ge 1$ there are odd positive $n_1 \equiv n_2 \pmod{2^m}$ whose first
division exponents are at least $m+1$ and exactly $m$. Hence no fixed residue class modulo $2^m$
determines the ledger (\Cref{thm:affine-ledger,thm:split-pair}).
\end{headline}

\begin{headline}[B: reverse-and-add in base 2]
Under binary reverse-and-add, the orbit of $22 = 10110_2$ reaches a four-word family after four
steps, cycles through its four phases with a parameter that grows by one per cycle, and never
reaches a palindrome (\Cref{thm:four-phase,cor:orbit-22}).
\end{headline}

\begin{headline}[C: Ducci nilpotency]
On cyclic binary vectors of length $2^k$, the map $x_i \mapsto x_i + x_{i+1}$ over
$\mathbb{F}_2$ vanishes after $2^k$ steps (\Cref{thm:ducci}).
\end{headline}

\begin{headline}[D: abelian move systems]
Consider a system of local moves in which any two distinct moves that are legal together remain
legal after each other and commute. Start from a configuration that admits no infinite legal run.
Then any two legal runs that end in stable configurations end in the same configuration, and they
fire every site the same number of times (\Cref{thm:G8}). Two small examples show that a unique
final configuration does not give canonical firing counts, and that termination does not give a
unique final configuration (\Cref{ex:G8-b,ex:G8-d}).
\end{headline}

\subsection{What is imported, and what is not claimed}\label{sec:intro-claims}

The deep theorems the catalog relies on are imported with attribution: Goodstein's theorem and the Kirby--Paris independence result, Dhar's abelian
property and the least-action principle of Fey, Levine and Peres, the solution of the angel
problem, the aperiodic tilings of Berger and Robinson, the de Bruijn--Erd\H{o}s compactness
theorem, and the density theorems of Bourgain and Kontorovich together with the reciprocity
obstruction of Haag, Kertzer, Rickards and Stange. The Collatz conjecture, the Erd\H{o}s--Straus
conjecture, the base-10 Lychrel problem, Recam\'an's coverage question, the Langton's-ant highway
conjecture, Gilbreath's conjecture and the chromatic number of the plane all remain open.

\paragraph{Roadmap.} \Cref{sec:related} places the paper in the literature, and \Cref{sec:setting}
fixes the vocabulary and gives the catalog. \Cref{sec:laws} treats the laws, \Cref{sec:theorems}
the general theorems and \Cref{sec:templates} the templates; the four headline results appear in
their entries. \Cref{sec:scope} lists inferences the results do not support,
and \Cref{sec:lean} records what is formalized. The supplement holds the source concordance, the
full register of non-implications, the case analyses and instance lists, the finite computations,
the reading of three laws as asymptotic statuses, and the replay commands.

\section{Related work}\label{sec:related}

\paragraph{The Six Birds series.} \FoundIV{} is the static catalog this paper extends. It presents
fifty-two laws in a six-part form --- setting, statement, proof, case analysis, relation to the six
structural roles, and limits --- each about fixed data \citep{Tsiokos2026FoundIV}. What separates each entry
here from \FoundIV{} is a quantifier over an unbounded run, scale or orbit. \FoundV{} studies
systems that carry out their own closure and repair \citep{FoundationsVCognition}; the laws here are
certified from outside, by an analyst auditing a run, and no law credits a system with certifying
itself. The quotient vocabulary of \Cref{sec:setting}, including split pairs and future-sufficient
quotients, is taken from \FoundIV{} and from the paper on holonomy with memory
\citep{Tsiokos2026Holonomy,Tsiokos2026Hiddenness}. Currencies as constraints, ledgers, potentials
and shadow prices come from \citet{Tsiokos2026Currency}, obstructions that survive every native
probe from \citet{Tsiokos2026NeedleKiller}, and the asymptotic statuses used in \Cref{supp:bridge}
from the paper on audited operational realisability \citep{Tsiokos2026AOR}.

\paragraph{Classical sources.} Each law with a classical core imports it. Termination by a
decreasing assignment into a well-founded order is Floyd's method \citep{Floyd1967}; Goodstein's
theorem and its unprovability in Peano arithmetic are due to \citet{Goodstein1944} and
\citet{KirbyParis1982}. The potential method of amortized analysis is Tarjan's
\citep{Tarjan1985,CormenEtAl}. The greedy expansion into unit fractions goes back to Fibonacci and
\citet{Sylvester1880}. The base-2 reverse-and-add result is recorded in the OEIS and on MathPages
\citep{OEIS-A060382,Brown-MathPages}; we know of no refereed proof, and the one in
\Cref{sec:reverse-add} is self-contained. The angel problem was settled by
\citet{Mathe2007}, \citet{Kloster2007}, \citet{Bowditch2007} and \citet{Gacs2007}. Order
independence for sandpiles is Dhar's abelian property \citep{Dhar1990}, and the least-action
principle is due to \citet{FeyLevinePeres2010}; abelian networks generalize both
\citep{HolroydEtAl2008}. The traffic automaton rule~184 is analyzed by \citet{Fuks1997}. Ducci
sequences are treated by \citet{ChamberlandThomas2004} and \citet{Breuer2007b}. Aperiodicity by
hierarchical forcing is due to \citet{Berger1966} and \citet{Robinson1971}. The de
Bruijn--Erd\H{o}s theorem is from \citet{BruijnErdos1951}, and the lower bound $5$ for the chromatic
number of the plane from \citet{DeGrey2018}. Thin orbits are treated by \citet{BourgainFuchs2010} and
\citet{BourgainKontorovich2014Apollonian}, and the reciprocity obstruction by \citet{HaagEtAl2024}.
For the Collatz problem we rely on the surveys of \citet{Lagarias1985,Lagarias2010} and on
\citet{Tao2019Collatz}.

\paragraph{Sources of the laws.} Each law isolates the structure behind a problem with a small rule
and a global obstruction: a tempting argument that relies on an auxiliary object, and the
certificate in the original system that would make the argument valid. \Cref{supp:method} works
this out for Collatz (G1) and sandpiles (G8) and records two related patterns that are not laws of
this paper.

\section{Setting}\label{sec:setting}

\paragraph{Runs.} A \emph{run} on a carrier $X$ is a sequence $x_0, x_1, x_2, \ldots$ with
$x_{k+1} = T(x_k)$ for an update $T : X \to X$, or $x_{k+1} = T_k(x_k)$ when the update may change
with the step. A set $A \subseteq X$ of terminal states is declared when the question is whether a
run stops. The laws of this paper quantify over whole runs: over every step, every legal order of
moves, every scale, or every orbit. That unbounded quantifier is what separates each law from its
nearest relative in \FoundIV{}, whose laws are about fixed data.

\paragraph{Quotients, descent and split pairs.} A \emph{quotient} is a map $q : X \to Q$ that
records what an observer can see. A map $F$ observed through a readout $r$ \emph{descends} through
$q$ when $r \circ F$ is a function of $q$, that is, when there is no \emph{split pair}: no $x, x'$
with $q(x) = q(x')$ and $r(F(x)) \ne r(F(x'))$ \citep{Tsiokos2026FoundIV}. When a whole family of
future readouts $r_j \circ F^{\,j}$ is considered, $q$ is \emph{future-sufficient} if every one of
them descends. Several laws are step-indexed versions of this test: a fixed quotient that fails to
be future-sufficient is exhibited by a split pair at some future step.

\paragraph{Hard hypotheses.} Each entry G1--G13 has a general statement. Its \emph{hard
hypothesis} is the hypothesis that carries the difficult content in applications --- a coverage
property, a closed family of states, a hierarchy of forced scales --- as opposed to routine data
such as a map or an order.

\paragraph{Laws, general theorems and templates.} Each entry receives one of three statuses,
decided by the mathematics.
\begin{itemize}
  \item A \emph{law} is a general statement together with a concrete system in which its hard
    hypothesis is proved in this paper, so that the conclusion is established for that system.
  \item A \emph{general theorem} is a statement whose hypotheses are ordinary data that can be
    checked system by system; its classical instances rest on imported results.
  \item A \emph{template} is an implication $H \Rightarrow C$ whose hard hypothesis $H$ is not
    established for any system of interest. $H$ is the \emph{open obligation} of the template. A
    template says nothing about a particular system until $H$ has been established for it.
\end{itemize}
A statement is a \emph{reformulation} when it is an equivalence obtained by unfolding definitions.
Which statements are also checked in Lean is a separate question, recorded in the last column of
\Cref{tab:catalog} and in \Cref{sec:lean}.

\paragraph{Instances.} An instance of an entry is a concrete system with its data. An instance is
proved here, imported from a cited theorem, or open.

\paragraph{Certificates.} Several entries concern a statement about an extended or auxiliary object
--- the $2$-adic integers, an ordinal, a potential --- that is used to settle a question about the
original system. What may be concluded about the original system is a \emph{certificate}: a
checkable statement in its own terms, such as an inequality between integers or the fact that a
run has stopped. The auxiliary object itself is never treated as a state of the original system.

\paragraph{The catalog at a glance.} \Cref{tab:catalog} lists the thirteen entries in the order in
which they are treated. The first column names each entry with a one-line gloss and the section
where it is treated. The second gives its status. The third names the instance proved in this
paper for a law, and the open obligation of a template, with the principal instances otherwise.
The fourth says what is checked in Lean. The full instance lists and case analyses are in
\Cref{supp:laws}.

{\footnotesize
\begin{longtable}{@{}>{\raggedright\arraybackslash}p{3.9cm}>{\raggedright\arraybackslash}p{1.3cm}>{\raggedright\arraybackslash}p{4.6cm}>{\raggedright\arraybackslash}p{4.4cm}@{}}
\caption{The catalog.}\label{tab:catalog}\\
\toprule
\textbf{Entry} & \textbf{Status} & \textbf{Instance or open obligation} & \textbf{Checked in Lean} \\
\midrule
\endfirsthead
\toprule
\textbf{Entry} & \textbf{Status} & \textbf{Instance or open obligation} & \textbf{Checked in Lean} \\
\midrule
\endhead
\bottomrule
\endfoot
\multicolumn{4}{@{}l}{\emph{Laws} (\Cref{sec:laws})} \\*
G3 \lawGThree: a potential bounds total cost over every finite stretch (\Cref{sec:G3}) & law &
binary counter (\Cref{ex:G3-counter}); doubling array & telescoping identity and bound over
integer costs; binary counter \\
G4 \lawGFour: covers that change with the target close an infinite family (\Cref{sec:G4}) & law &
initial intervals (\Cref{ex:G4-intervals}); greedy unit fractions; Erd\H{o}s--Straus open & general
statement; initial intervals \\
G5 \lawGFive: a closed family of states that avoids the target confines the run (\Cref{sec:G5}) &
law & orbit of $22$ in base $2$ (\Cref{cor:orbit-22}); $196$ in base $10$ open; no-go template &
confinement; four-phase family; orbit of $22$; no-go template \\
G8 \lawGEight: commuting moves give an order-independent end state and firing counts
(\Cref{sec:G8}) & law & sandpiles (\Cref{ex:G8-sandpile}) & order independence; the two separating
examples; a three-site sandpile; not general graphs \\
G10 \lawGTen: a relation survives every iterate of a lossy map if an invariant is regenerated
(\Cref{sec:G10}) & law & checksum cascade (\Cref{ex:G10-checksum}); Ducci collapse; Gilbreath open &
general statement; checksum cascade; Ducci over $\mathbb{F}_2$; integer Ducci for $n = 4$ \\
G12 \lawGTwelve: a finite obstruction forbids a global colouring, and so do its images
(\Cref{sec:G12}) & law & Moser spindle, $\chi(\mathbb{R}^2) \ge 4$ (\Cref{ex:G12-moser});
$\chi(\mathbb{R}^2) \ge 5$ imported & general statement and the equivalence given compactness; the
spindle graph has no $3$-colouring; not the plane embedding \\
G13 \lawGThirteen: a sparse orbit misses only a density-zero part of a dense target
(\Cref{sec:G13}) & law, weak instance & powers of two, an arithmetic example
(\Cref{ex:G13-powers}); thin-orbit saturation a template & density lemma; powers of two \\
\midrule
\multicolumn{4}{@{}l}{\emph{General theorems} (\Cref{sec:theorems})} \\*
G2 \lawGTwo: a value in a well-founded order that falls at every step forces termination
(\Cref{sec:G2}) & theorem & Goodstein sequences and Kirby--Paris (imported) & general statement,
without extracting the terminal stage \\
G7 \lawGSeven: escape and confinement strategies against an adversary (\Cref{sec:G7}) & theorem &
the angel problem (imported) & general statement \\
\midrule
\multicolumn{4}{@{}l}{\emph{Templates} (\Cref{sec:templates})} \\*
G1 \lawGOne: every nonterminal point eventually earns a certificate of descent (\Cref{sec:G1}) &
template & (H1) and (H2) of \Cref{thm:G1b}; for Collatz both open & reformulation; template; Collatz
ledger and split pairs \\
G6 \lawGSix: obstructions made by the run's own history are covered or leave holes
(\Cref{sec:G6}) & template & bridge from rate conditions to the premises of \Cref{thm:G6};
Recam\'an open & the two final steps \\
G9 \lawGNine: local defects resolve or turn into certified motion (\Cref{sec:G9}) & template & (H1)
of \Cref{thm:G9}; evacuation of every low-density rule-184 ring & census bound; packaging step;
rule-184 partial results, car conservation, evacuation for $n \le 8$ \\
G11 \lawGEleven: local rules with a forced hierarchy exclude every period (\Cref{sec:G11}) &
template & a hierarchy for a given tile set; Robinson and others (imported) & the implication \\
\end{longtable}
}

\paragraph{Map of results.} \Cref{tab:map} lists the numbered results. Each is proved where it is stated.

{\footnotesize
\begin{longtable}{@{}>{\raggedright\arraybackslash}p{3.2cm}>{\raggedright\arraybackslash}p{11.3cm}@{}}
\caption{Map of results.}\label{tab:map}\\
\toprule
\textbf{Result} & \textbf{What it establishes} \\
\midrule
\endfirsthead
\toprule
\textbf{Result} & \textbf{What it establishes} \\
\midrule
\endhead
\bottomrule
\endfoot
\Cref{thm:G3} & Total cost over a finite stretch equals amortized cost plus the drop in potential, so amortized costs at most $B$ give total cost at most $nB+\Phi(s_0)$. \\
\Cref{thm:G4} & Patches that prove the target on their regions and cover every target prove it everywhere; under (H4) no finite list from the restricted class covers. \\
\Cref{thm:G5} & A run that enters a closed pattern containing no target stays in it and never reaches a target. \\
\Cref{thm:four-phase} & Binary reverse-and-add maps the four words $P_0(r),\dots,P_3(r)$ cyclically into the next family, and none is a palindrome. \\
\Cref{cor:orbit-22} & The orbit of $22$ in base $2$ enters the four-phase family and never reaches a palindrome. \\
\Cref{thm:G8} & If legal moves commute and no infinite legal run exists, all stabilizing runs end in the same configuration with the same counter vector. \\
\Cref{thm:G10} & A relation carried through one lossy step by a regenerated invariant persists through every iterate. \\
\Cref{thm:ducci} & Over $\mathbb F_2$ the Ducci map is nilpotent on words of length $2^k$. \\
\Cref{thm:G12} & A finite witness rules out a $k$-colouring, and its image under any map preserving $E$ is again a witness. \\
\Cref{prop:G12-compact} & Under compactness, no $k$-colouring exists exactly when a finite witness does. \\
\Cref{thm:G13-density} & A density-zero set is density zero relative to any target of positive lower density. \\
\Cref{thm:G2} & A value in a well-founded order that falls at every step outside the terminal set forces termination. \\
\Cref{thm:G7} & An invariant set closed under the angel's strategy and every legal devil reply certifies escape; the dual certifies confinement. \\
\Cref{thm:affine-ledger} & After $k$ odd steps of the accelerated Collatz map, $2^{A_k}n_k=3^kn+B_k$. \\
\Cref{thm:split-pair} & For each $m$, two odd numbers congruent modulo $2^m$ have different division exponents. \\
\Cref{cor:no-residue} & The division exponent is not a function of $n$ modulo $2^m$ for any $m$. \\
\Cref{prop:G1} & The system is live exactly when no point lies in every bad tail $N_k$. \\
\Cref{thm:G6} & Holes that are eventually reduced are eventually removed, and an unvisited point that stays unvisited is never visited. \\
\Cref{thm:G9-census} & A nonincreasing count of unresolved defects is bounded. \\
\Cref{prop:rule184} & Rule 184 keeps jam-free rings jam-free and conserves cars; on rings of $2$ to $8$ cells, every configuration with at most half the cells occupied is jam-free after $n$ steps. \\
\end{longtable}
}


\section{Laws with instances}\label{sec:laws}

Each of the seven entries in this section has a general statement and a concrete system in which
its hard hypothesis is proved, in the text below. Hypotheses are labelled (H1), (H2), \dots\ within
each entry.

\subsection{G3: \lawGThree}\label{sec:G3}

\emph{Name.} A potential acts as a \emph{currency}: credit stored in the state pays for expensive
steps later.

Let $s_0 \to s_1 \to \cdots \to s_n$ be a finite stretch of a run in which step $i$ has cost
$a_i$, and let $\Phi \ge 0$ be a potential on states. The \emph{amortized cost} of step $i$ is
$\widehat a_i = a_i + \Phi(s_i) - \Phi(s_{i-1})$.

\begin{theorem}[Amortized potential]\label{thm:G3}
For every finite stretch,
\[
  \sum_{i=1}^{n} a_i \;=\; \sum_{i=1}^{n} \widehat a_i + \Phi(s_0) - \Phi(s_n)
  \;\le\; \sum_{i=1}^{n} \widehat a_i + \Phi(s_0).
\]
If \textup{(H1)} $\widehat a_i \le B$ at every step, then $\sum_{i=1}^n a_i \le nB + \Phi(s_0)$.
\end{theorem}

\begin{proof}
The potential terms in $\sum \widehat a_i$ telescope, and $\Phi(s_n) \ge 0$.
\end{proof}

This is the potential method of amortized analysis \citep{Tarjan1985,CormenEtAl}. G3 is a safety
statement about every finite stretch and says nothing about eventual success; it assumes the
potential rather than constructing it. A potential counts as a certificate only when it is
nonnegative, enters through the exact telescoping formula, and bounds the cost claim actually made.
A statistic that merely correlates with cost does not.

\begin{example}[Binary counter]\label{ex:G3-counter}
Let the state be a binary counter, let an increment cost the number of bits it changes, and let
$\Phi$ be the number of $1$ bits. An increment that turns $r$ trailing ones into zeros and one zero
into a one costs $r + 1$ and changes $\Phi$ by $1 - r$, so its amortized cost is exactly $2$ and
\textup{(H1)} holds with $B = 2$. Hence $n$ increments from zero change at most $2n$ bits, counting a
newly created leading bit as a changed bit.
\end{example}

\Cref{fig:counter} shows the first eight increments.

\begin{figure}[htbp]
\centering
\begin{tikzpicture}
  \node[inner sep=0pt] {\footnotesize
    \begin{tabular}{@{}c c c c c c c@{}}
      \toprule
      $n$ & counter $n-1\to n$ & bit positions flipped & cost $a_n$ & $\Phi$ before $\to$ after & $\Delta\Phi$ & $a_n+\Delta\Phi$ \\
      \midrule
      $1$ & $0000\to0001$ & $\{0\}$       & $1$ & $0\to1$ & $+1$ & $2$ \\
      $2$ & $0001\to0010$ & $\{0,1\}$     & $2$ & $1\to1$ & $0$  & $2$ \\
      $3$ & $0010\to0011$ & $\{0\}$       & $1$ & $1\to2$ & $+1$ & $2$ \\
      $4$ & $0011\to0100$ & $\{0,1,2\}$   & $3$ & $2\to1$ & $-1$ & $2$ \\
      $5$ & $0100\to0101$ & $\{0\}$       & $1$ & $1\to2$ & $+1$ & $2$ \\
      $6$ & $0101\to0110$ & $\{0,1\}$     & $2$ & $2\to2$ & $0$  & $2$ \\
      $7$ & $0110\to0111$ & $\{0\}$       & $1$ & $2\to3$ & $+1$ & $2$ \\
      $8$ & $0111\to1000$ & $\{0,1,2,3\}$ & $4$ & $3\to1$ & $-2$ & $2$ \\
      \bottomrule
    \end{tabular}};
\end{tikzpicture}
\caption{Eight increments from zero, with bit positions numbered from the least significant bit. The displayed potential is the number of $1$ bits.}
\label{fig:counter}
\end{figure}


For a doubling array with $\Phi = \max(0, 2\,\mathrm{size} - \mathrm{capacity})$, every insertion
has amortized cost at most $3$ by the same kind of computation \citep{CormenEtAl}.

\subsection{G4: \lawGFour}\label{sec:G4}

\emph{Name.} The cover of the target family must \emph{move}: which patch covers a target depends
on the target.

Let $\{P(t) : t \in \mathcal T\}$ be an infinite family of targets. A \emph{patch} $\alpha$ comes
with a region $R_\alpha \subseteq \mathcal T$ on which it claims to prove $P$.

\begin{theorem}[Moving-cover exhaustion]\label{thm:G4}
Consider the hypotheses
\begin{itemize}
  \item \textup{(H1)} every patch proves $P(t)$ for every $t$ in its region;
  \item \textup{(H2)} every target lies in the region of some patch;
  \item \textup{(H3)} if $P$ holds everywhere, the declared patches cover every target;
  \item \textup{(H4)} for a declared restricted class of patches, every finite list from the class
    misses some target.
\end{itemize}
Then: (H1) and (H2) give $P(t)$ for every $t$; under (H1) and (H3), $P$ holds everywhere exactly
when the patches cover every target; and under (H4) no finite list from the restricted class covers
every target.
\end{theorem}

\begin{proof}
The first claim applies (H1) to the covering patch. The forward direction of the second is (H3),
and its converse is the first claim. The third is (H4) applied to the given list.
\end{proof}

Each conclusion is close to its hypothesis: (H2) is total coverage, and (H4) already says that every
finite list misses a target. The substance lies in establishing (H2) and (H4) for a system, usually
by a residual for each target that shrinks as the cover grows.

\begin{example}[Initial intervals]\label{ex:G4-intervals}
Let $\mathcal T = \mathbb{N}$, and let the patch with index $a$ cover the interval $[0, a)$. For a
target $t$ the residual $\max(0, t + 1 - a)$ decreases strictly as $a$ grows while it is positive,
and vanishes at $a = t + 1$, so (H2) holds. A finite list of patches misses the largest of its
endpoints, so (H4) holds for the class of all patches. Initial intervals therefore exhaust the
natural numbers as their endpoints grow, although no finite collection of them covers every natural
number.
\end{example}

The greedy expansion of a fraction $p/q \in (0,1)$ into distinct unit fractions is the classical
instance. With $m = \lceil q/p \rceil$ the remainder after subtracting $1/m$ is $(pm - q)/(qm)$,
and $0 \le pm - q < p$, so the numerator strictly decreases and the expansion ends
\citep{Sylvester1880}. A fixed finite set of denominators has finitely many subset sums, so no fixed
package covers all targets; the greedy cover succeeds because its denominators depend on the
target. For the Erd\H{o}s--Straus conjecture ($4/n = 1/x + 1/y + 1/z$ for all $n \ge 2$), a
solution for a prime yields one for its multiples, and no finite family of polynomial identities can
cover the class $n \equiv 1$ modulo every prime because $1$ is a quadratic residue
\citep{Mordell1969}. The conjecture is open, and density results
\citep{Vaughan1970,ElsholtzTao2013} do not give it.

\subsection{G5: \lawGFive}\label{sec:G5}

\emph{Name.} Predicting digit dynamics can need digits from ever deeper positions, as a carry
travels; the \emph{horizon} of useful digits recedes. A closed family of states \emph{confines} the
run when prediction fails.

Let $R : X \to X$ be an update and let a set of \emph{targets} be declared; for reverse-and-add the
targets are the palindromes.

\begin{theorem}[Confinement by a closed pattern]\label{thm:G5}
Let $\mathcal P \subseteq X$ satisfy \textup{(H1)} $R(\mathcal P) \subseteq \mathcal P$ and
\textup{(H2)} no element of $\mathcal P$ is a target. If $R^N(x) \in \mathcal P$, then
$R^{N+t}(x) \in \mathcal P$ and $R^{N+t}(x)$ is not a target for every $t \ge 0$.
\end{theorem}

\begin{proof}
Induction on $t$.
\end{proof}

\subsubsection*{The instance: reverse-and-add in base 2}\label{sec:reverse-add}

A \emph{binary word} here is a word over $\{0,1\}$ that begins with $1$; every positive integer has
exactly one. For a word $w$ let $\bar w$ be its reversal, which may begin with zeros. Binary
reverse-and-add sends $w$ to the binary word of $w + \bar w$, reading both as numbers; we write
$R$ for this map on words and on the positive integers they represent. A word is a
\emph{palindrome} when it equals its reversal. For $r \ge 2$ define the four words
\[
\begin{aligned}
  P_0(r) &= 10\,1^r\,01\,0^r, &\qquad
  P_1(r) &= 11\,0^{r-2}\,1\,000\,1^{r-2}\,01, \\
  P_2(r) &= 10\,1^r\,01\,0^{r+1}, &
  P_3(r) &= 11\,0^r\,10\,1^{r-1}\,01,
\end{aligned}
\]
where $b^j$ denotes $j$ copies of the digit $b$.

\begin{theorem}[Four-phase family]\label{thm:four-phase}
For every $r \ge 2$, as binary words,
\[
  R(P_0(r)) = P_1(r), \qquad R(P_1(r)) = P_2(r), \qquad R(P_2(r)) = P_3(r), \qquad
  R(P_3(r)) = P_0(r+1),
\]
and none of the four words is a palindrome.
\end{theorem}

\begin{proof}
Put $x = 2^r$. The value of each word and of its reversal is a polynomial in $x$ with integer
coefficients, given in \Cref{tab:four-phase}. In each row the value and the reversal add up to the
value of the next word; for the last row, $P_0(r+1)$ has value $12(2x)^2 - 3(2x) = 48x^2 - 6x$.
Since a binary word is determined by its value, each identity of values is an identity of words.
Finally, $P_0$ and $P_2$ begin with $1$ and end with $0$, and $P_1$ and $P_3$ begin with $11$ and
end with $01$, so none equals its reversal.
\end{proof}

\begin{table}[htbp]
\centering
\small
\caption{Values of the four phases and their reversals, with $x = 2^r$.}\label{tab:four-phase}
\begin{tabular}{@{}lllll@{}}
\toprule
Word & Value & Reversal & Sum & Next word \\
\midrule
$P_0(r)$ & $12x^2 - 3x$ & $12x - 3$ & $12x^2 + 9x - 3$ & $P_1(r)$ \\
$P_1(r)$ & $12x^2 + 9x - 3$ & $12x^2 - 15x + 3$ & $24x^2 - 6x$ & $P_2(r)$ \\
$P_2(r)$ & $24x^2 - 6x$ & $12x - 3$ & $24x^2 + 6x - 3$ & $P_3(r)$ \\
$P_3(r)$ & $24x^2 + 6x - 3$ & $24x^2 - 12x + 3$ & $48x^2 - 6x$ & $P_0(r+1)$ \\
\bottomrule
\end{tabular}
\end{table}

The family is recorded in the OEIS and on MathPages \citep{OEIS-A060382,Brown-MathPages}; the proof
here is self-contained.

\begin{corollary}[The orbit of $22$]\label{cor:orbit-22}
The orbit of $22$ under binary reverse-and-add begins
\[
  10110 \;\to\; 100011 \;\to\; 1010100 \;\to\; 1101001 \;\to\; 10110100 = P_0(2),
\]
thereafter runs through $P_1(2), P_2(2), P_3(2), P_0(3), \ldots$, and never reaches a palindrome.
\end{corollary}

\begin{proof}
The first four steps are direct computations ($22 + 13 = 35$, $35 + 49 = 84$, $84 + 21 = 105$,
$105 + 75 = 180$), and none of the four initial words is a palindrome. From $P_0(2)$ on, the family
of \Cref{thm:four-phase} satisfies (H1) and (H2) of \Cref{thm:G5}, so no later word is a
palindrome.
\end{proof}

In base $10$ no such certificate is known: whether the orbit of $196$ reaches a palindrome is
open, and finite searches are evidence only \citep{Walker-Fourmilab}.

\subsubsection*{Why a certificate of this kind is needed}

Let $q_m$ be a family of quotients --- for reverse-and-add in base $b$, $q_m(n) = n \bmod b^m$
records the last $m$ digits --- and let $r_j$ be a family of future readouts, such as the last
digits after $j$ steps.

\begin{schema}[No fixed prediction depth]\label{thm:G5-nogo}
Assume
\begin{itemize}
  \item \textup{(H3)} for every depth $m$ there are a step $j$ and states $y, z$ with
    $q_m(y) = q_m(z)$ and $r_j(y) \ne r_j(z)$.
\end{itemize}
Then no depth $m$ is future-sufficient for the declared readouts.
\end{schema}

\begin{proof}
The pair $y, z$ at depth $m$ is a split pair for $q_m$.
\end{proof}

The template holds for any quotients and readouts. For the Collatz exponent, with
$q_m(n) = n \bmod 2^m$ and the readout $a(n)$, (H3) is \Cref{thm:split-pair}. Whether (H3) holds
for base-$2$ reverse-and-add with a stated family of readouts is open; the confinement of $22$ does
not need it.

\subsection{G8: \lawGEight}\label{sec:G8}\label{sec:abelian}

\emph{Name.} The \emph{odometer} is the vector of firing counts; \emph{abelianization} means that
the order of commuting moves stops mattering.

\begin{definition}[Move system]\label{def:move-system}
A \emph{move system} consists of a set $C$ of configurations, a set $V$ of sites, and for each
$v \in V$ a move $m_v : C \to C$ with a legality predicate $L_v$ on $C$. A \emph{legal run} from $x$
is a finite sequence of sites $v_0, \ldots, v_{\ell-1}$ with $L_{v_i}$ holding at the configuration
reached before step $i$. Its \emph{counter vector} records how often each site occurs. A
configuration is \emph{stable} when no move is legal there.
\end{definition}

\begin{definition}[Abelian compatibility]\label{def:abelian}
The move system satisfies \textup{(H1)} when for every configuration $c$ and all \emph{distinct}
sites $v \ne w$ that are both legal at $c$: $w$ is legal at $m_v(c)$, $v$ is legal at $m_w(c)$,
and $m_w(m_v(c)) = m_v(m_w(c))$.
\end{definition}

\begin{theorem}[Order independence]\label{thm:G8}
Assume \textup{(H1)} and \textup{(H2)}: no infinite sequence of legal moves starts at $x$. If two
legal runs from $x$ end in stable configurations, then they end in the same configuration and have
the same counter vector.
\end{theorem}

\begin{proof}
Induction on the length of the first run. If it is empty, $x$ is stable and the second run is empty
as well. Otherwise let $v$ be its first site. The second run ends in a stable configuration, so $v$
is not legal at its end. By (H1), firing any site other than $v$ keeps $v$ legal, so $v$ occurs in
the second run. Moving that occurrence step by step to the front uses (H1) at each swap, keeps the
run legal, and changes neither its final configuration nor its counter vector. Both runs now begin
with $v$, and the induction hypothesis applies to their remainders from $m_v(x)$.
\end{proof}

This is the abstract form of Dhar's abelian property for sandpiles \citep{Dhar1990}. Under (H2)
every maximal legal run from $x$ is finite and ends in a stable configuration, so all of them end in
one configuration $x^\circ$ with one counter vector $u_x$, the \emph{odometer}.

\begin{example}[Sandpiles]\label{ex:G8-sandpile}
Let $G$ be a finite connected graph with a distinguished sink $s$. A configuration assigns a number
of chips $\eta(v) \ge 0$ to each $v \ne s$. A site $v$ is legal when $\eta(v) \ge \deg(v)$, and
toppling it sends one chip along each edge at $v$; chips sent to $s$ are lost. For distinct legal
$v, w$, toppling $v$ does not decrease $\eta(w)$, so $w$ stays legal, and both orders subtract the
same vector of chip changes; this is (H1). For (H2), suppose an infinite legal run exists. The
number of chips never increases, so every $\eta(v)$ stays bounded. A site that topples infinitely
often sends infinitely many chips to each neighbour, which must then topple infinitely often too.
By connectedness some neighbour of $s$ topples infinitely often, and each such toppling loses a chip
to $s$; this contradicts the finiteness of the chip count. \Cref{thm:G8} therefore gives every
sandpile a unique stable configuration and odometer.

One small case in which sites exchange chips has three mutually adjacent sites, each also
joined to the sink, so that every degree is $3$. Toppling a site removes three of its chips and adds
one to each other site, and one chip is lost to the sink; so every toppling lowers the total by
exactly one, which gives (H2) at once.
\end{example}

\Cref{fig:sandpile} shows one toppling in this case.

\begin{figure}[htbp]
\centering
\begin{tikzpicture}[>=Stealth, site/.style={circle,draw,fill=white,minimum size=8mm,inner sep=0pt,font=\small}]
  \begin{scope}
    \coordinate (a) at (0,1.65);
    \coordinate (b) at (-1.55,-0.85);
    \coordinate (c) at (1.55,-0.85);
    \coordinate (s) at (0,-0.05);
    \draw[gray] (a)--(b)--(c)--(a) (a)--(s) (b)--(s) (c)--(s);
    \draw[->,line width=0.7pt,shorten <=12pt,shorten >=12pt] (a)--(b);
    \draw[->,line width=0.7pt,shorten <=12pt,shorten >=12pt] (a)--(c);
    \draw[->,line width=0.7pt,shorten <=12pt,shorten >=12pt] (a)--(s);
    \node[site,fill=gray!25] at (a) {3};
    \node[site] at (b) {0};
    \node[site] at (c) {0};
    \node[site,fill=gray!12] at (s) {$s$};
    \node[font=\footnotesize,above=15pt] at (a) {$v_1$};
    \node[font=\footnotesize,below left=13pt] at (b) {$v_2$};
    \node[font=\footnotesize,below right=13pt] at (c) {$v_3$};
    \node[font=\small] at (0,3) {before};
  \end{scope}
  \draw[->,line width=0.7pt] (2.3,0.35)--(3.5,0.35)
    node[midway,above,font=\footnotesize] {topple $v_1$};
  \begin{scope}[xshift=5.8cm]
    \coordinate (aa) at (0,1.65);
    \coordinate (bb) at (-1.55,-0.85);
    \coordinate (cc) at (1.55,-0.85);
    \coordinate (ss) at (0,-0.05);
    \draw[gray] (aa)--(bb)--(cc)--(aa) (aa)--(ss) (bb)--(ss) (cc)--(ss);
    \node[site] at (aa) {0};
    \node[site,fill=gray!12] at (bb) {1};
    \node[site,fill=gray!12] at (cc) {1};
    \node[site,fill=gray!12] at (ss) {$s$};
    \node[font=\footnotesize,above=15pt] at (aa) {$v_1$};
    \node[font=\footnotesize,below left=13pt] at (bb) {$v_2$};
    \node[font=\footnotesize,below right=13pt] at (cc) {$v_3$};
    \node[font=\small] at (0,3) {after};
  \end{scope}
\end{tikzpicture}
\caption{One toppling on the three-site graph: $(3,0,0)\to(0,1,1)$. The arrows show one chip sent along each edge from $v_1$, including the edge to the sink $s$.}
\label{fig:sandpile}
\end{figure}


Rotor-routing belongs to the same family of abelian networks \citep{HolroydEtAl2008}. The next two
examples separate the theorem from weaker properties of rewriting systems.

\begin{example}[A unique final configuration without canonical counts]\label{ex:G8-b}
Take configurations $S, A, B, N$ and legal moves $r : S \to A$, $s : S \to B$, $p : A \to N$,
$t : A \to N$, $q : B \to N$, and no others. Every run terminates, and every maximal run ends at $N$.
The three maximal runs $rp$, $rt$ and $sq$ have three different counter vectors. (H1) fails at
$S$: after $r$ the move $s$ is no longer legal.
\end{example}

\begin{example}[Termination without a unique final configuration]\label{ex:G8-d}
Take configurations $S, A, B$ and legal moves $r : S \to A$ and $s : S \to B$ only. Both runs
terminate, at the distinct stable configurations $A$ and $B$.
\end{example}

\begin{remark}[Least action]\label{rem:least-action}
The least-action principle of \citet{FeyLevinePeres2010} compares the odometer with every
nonnegative script of firings that stabilizes the configuration, including scripts that pass
through illegal intermediate configurations. It is imported. The weaker comparison with scripts
that are themselves realized by legal stabilizing runs adds nothing under the hypotheses of
\Cref{thm:G8}, which make all such scripts equal.
\end{remark}

\subsection{G10: \lawGTen}\label{sec:G10}

\emph{Name.} A boundary relation \emph{persists} through a map that is \emph{lossful}, that is,
not injective.

Let $D : X \to X$ be a map, $b : X \to B$ a boundary readout, $\rho : B \to B$ a boundary update, and
$I$ a predicate on $X$, the interior invariant. The run is $x_{n+1} = D(x_n)$.

\begin{theorem}[Persistence through iterated loss]\label{thm:G10}
Assume
\begin{itemize}
  \item \textup{(H1)} $I(x)$ implies $b(D(x)) = \rho(b(x))$, and
  \item \textup{(H2)} $I(x)$ implies $I(D(x))$,
\end{itemize}
and $I(x_0)$. Then $I(x_n)$ and $b(x_n) = \rho^n(b(x_0))$ for every $n \ge 0$.
\end{theorem}

\begin{proof}
Induction on $n$.
\end{proof}

The theorem holds for every map $D$. It is of interest when $D$ is not injective, so that interior
information is lost while the boundary relationship survives because the invariant is regenerated
at every step.

\begin{example}[Checksum cascade]\label{ex:G10-checksum}
On $X = \{0,1\}^2 \times \{0,1\}$ let $D((a, b), c) = ((c, 0), c)$, with boundary $b((a,b),c) = c$,
$\rho$ the identity and invariant $c = a \oplus b$. The map is not injective, and (H1) and (H2)
hold: if $c = a \oplus b$ then the new state $((c, 0), c)$ satisfies $c = c \oplus 0$. The checksum
therefore persists forever while the two data bits are lost.
\end{example}

\subsubsection*{A collapse example: Ducci maps}\label{sec:ducci}

For $n \ge 1$ let $D$ act on cyclic binary vectors $x = (x_0, \ldots, x_{n-1}) \in
\mathbb{F}_2^{\,n}$ by $D(x)_i = x_i + x_{i+1}$, indices modulo $n$.

\begin{theorem}[Ducci nilpotency over $\mathbb{F}_2$]\label{thm:ducci}
If $n = 2^k$, then $D^{\,n}(x) = 0$ for every $x \in \mathbb{F}_2^{\,n}$.
\end{theorem}

\begin{proof}
Write $D = I + S$ with $S$ the cyclic shift. Over $\mathbb{F}_2$, induction on $j$ gives
$D^{2^j}(x)_i = x_i + x_{i+2^j}$, since squaring $I + S^{2^j}$ gives $I + S^{2^{j+1}}$. For $j = k$
the shift $S^{2^k} = S^n$ is the identity, so $D^n(x)_i = 2x_i = 0$.
\end{proof}

The result is classical \citep{ChamberlandThomas2004,Breuer2007b}. For vectors of length $2^k$ the
parity structure is destroyed completely, so no nonzero persistence claim for all such vectors can
rest on parity.

\begin{remark}[Integer Ducci tuples]\label{rem:ducci-integer}
The integer Ducci map sends $(x_0, \ldots, x_{n-1})$ to $(|x_0 - x_1|, \ldots, |x_{n-1} - x_0|)$.
Modulo $2$ it is $D$, and it commutes with multiplication by $2$. For $n = 2^k$, \Cref{thm:ducci}
makes every entry even after $n$ steps, after which the run is twice another Ducci run. Since the
entries never exceed the largest entry of the first iterate, they become divisible by arbitrarily
large powers of $2$ and hence vanish. For $n = 4$ the argument gives an explicit bound: every
$x \in \mathbb{N}^4$ reaches the zero tuple after $4m$ steps of the integer map, where $m$ is the
largest entry of $x$, because after
$4k$ steps every entry is a multiple of $2^k$ and at most $m$, and $2^m > m$. For other lengths
integer tuples can cycle without reaching zero, for example $(0,1,1) \to (1,0,1) \to (1,1,0) \to
(0,1,1)$.
\end{remark}

Gilbreath's conjecture, read as a persistence claim for the leading entry of iterated absolute
differences of the primes, is open \citep{Odlyzko1993}.

\subsection{G12: \lawGTwelve}\label{sec:G12}

\emph{Name.} A finite witness \emph{radiates}: its images under the symmetries of the graph are
witnesses too.

Let $E$ be a graph relation on a set $X$. A \emph{proper $k$-colouring} is a map $c : X \to \{1,
\ldots, k\}$ with $c(x) \ne c(y)$ whenever $E(x, y)$. A \emph{finite witness} for $k$ is a finite set
$W \subseteq X$ whose induced graph has no proper $k$-colouring. A map $\gamma : X \to X$
\emph{preserves} $E$ when $E(x, y) \Leftrightarrow E(\gamma x, \gamma y)$.

\begin{theorem}[Finite witnesses]\label{thm:G12}
Let $W$ be a finite witness for $k$.
\begin{enumerate}
  \item $X$ has no proper $k$-colouring.
  \item If $\gamma$ preserves $E$, then $\gamma(W)$ is a finite witness for $k$.
\end{enumerate}
\end{theorem}

\begin{proof}
A proper colouring of $X$ restricts to one of $W$. A proper colouring $c$ of $\gamma(W)$ pulls back
to the proper colouring $c \circ \gamma$ of $W$, because $\gamma$ sends edges to edges; $\gamma$
need not be injective.
\end{proof}

\begin{example}[The Moser spindle]\label{ex:G12-moser}
In the plane with $E(x, y)$ meaning $\lVert x - y \rVert = 1$, take two rhombi with vertices $A, B, D, C$ and
$A, B', D', C'$ in cyclic order and unit sides, each made of two equilateral unit triangles ($ABC$,
$BCD$ and $AB'C'$, $B'C'D'$), so that $A$ and $D$ are opposite vertices with
$\lVert A - D \rVert = \sqrt 3$, and likewise $\lVert A - D' \rVert = \sqrt 3$. Rotate the second
rhombus about $A$ by the angle $\theta$ with $2\sqrt 3 \sin(\theta/2) = 1$, which exists because
$1 < 2\sqrt 3$; then $\lVert D - D' \rVert = 1$. The seven points $A, B, C, D, B', C', D'$ are
distinct, and the eleven pairs $AB, AC, BC, BD, CD$, $AB', AC', B'C', B'D', C'D'$ and $DD'$ are at
unit distance. In a proper
$3$-colouring, $B$ and $C$ take the two colours other than that of $A$, so $D$ takes the colour of
$A$; likewise $D'$ does, and $D, D'$ are adjacent. So the seven points form a finite witness for
$k = 3$, and $\chi(\mathbb{R}^2) \ge 4$ by \Cref{thm:G12}; by part~2, so does every rotated and
translated copy.
\end{example}

\Cref{fig:moser} shows the seven points and the eleven unit distances.

\begin{figure}[htbp]
\centering
\begin{tikzpicture}[scale=3.3, every node/.style={font=\small}]
  % The long diagonals AD and AD' have length sqrt(3). Their angles are
  % -alpha and +alpha, where alpha = asin(1/(2 sqrt(3))).
  \coordinate (A)  at (0,0);
  \coordinate (B)  at (0.973493765,0.228713554);
  \coordinate (C)  at (0.684818630,-0.728713554);
  \coordinate (D)  at (1.658312395,-0.500000000);
  \coordinate (Bp) at (0.684818630,0.728713554);
  \coordinate (Cp) at (0.973493765,-0.228713554);
  \coordinate (Dp) at (1.658312395,0.500000000);
  \draw[line width=0.5pt] (A)--(B)--(D)--(C)--cycle;
  \draw[line width=0.5pt] (B)--(C);
  \draw[line width=0.5pt] (A)--(Bp)--(Dp)--(Cp)--cycle;
  \draw[line width=0.5pt] (Bp)--(Cp);
  \draw[line width=1pt] (D)--(Dp);
  \fill (A) circle (0.018) node[left=3pt] {$A$};
  \fill (B) circle (0.018) node[above right=2pt] {$B$};
  \fill (C) circle (0.018) node[below left=2pt] {$C$};
  \fill (D) circle (0.018) node[below right=2pt] {$D$};
  \fill (Bp) circle (0.018) node[above left=2pt] {$B'$};
  \fill (Cp) circle (0.018) node[below right=2pt] {$C'$};
  \fill (Dp) circle (0.018) node[above right=2pt] {$D'$};
\end{tikzpicture}
\caption{Two unit rhombi sharing $A$, with the second rotated by $\theta$ where $2\sqrt3\sin(\theta/2)=1$. The heavier edge joins $D$ and $D'$; all eleven drawn edges have unit length.}
\label{fig:moser}
\end{figure}


\begin{proposition}[Completeness of finite witnesses under compactness]\label{prop:G12-compact}
Assume
\begin{itemize}
  \item \textup{(H1)} if $X$ has no proper $k$-colouring, then some finite witness for $k$ exists.
\end{itemize}
Then $X$ has no proper $k$-colouring exactly when a finite witness for $k$ exists.
\end{proposition}

\begin{proof}
The forward direction is the hypothesis and the backward direction is \Cref{thm:G12}.
\end{proof}

The hypothesis is the de~Bruijn--Erd\H{o}s theorem \citep{BruijnErdos1951}, which uses the axiom of
choice. For the unit-distance graph of the plane the current bounds are $5 \le \chi(\mathbb{R}^2)
\le 7$. The lower bound is a finite witness in exactly the sense of \Cref{thm:G12}: de~Grey
constructed a $1581$-vertex unit-distance graph that is not $4$-colourable \citep{DeGrey2018}, and
later work reduced this to $553$ vertices \citep{Heule2018} and to $509$ \citep{Parts2020}. The
upper bound is the classical hexagonal colouring. The exact value is open, and questions of this
kind can depend on the set-theoretic background \citep{ShelahSoifer2003,Falconer1981}; the finite
lower bound does not, once the graph and its non-colourability are accepted.

\subsection{G13: \lawGThirteen}\label{sec:G13}

\emph{Name.} An orbit that is \emph{thin} (sparse in a precise sense) can still \emph{saturate} a
target, missing only a vanishing proportion of it.

Let $m, t : \mathbb{N} \to \mathbb{N}$ count, up to $N$, the missed members and all members of a
target set.

\begin{theorem}[Relative density]\label{thm:G13-density}
If $m(N) = o(N)$ and $t(N) \ge cN$ for all $N$ with a fixed $c > 0$, then $m(N) = o(t(N))$.
\end{theorem}

\begin{proof}
Given $\varepsilon > 0$, eventually $m(N) \le \varepsilon c N \le \varepsilon\, t(N)$.
\end{proof}

\begin{example}[Powers of two]\label{ex:G13-powers}
Let the target be the positive integers and let the represented set be the integers that are not
powers of two. The missed family, the powers of two, is infinite, and its count up to $N$ is at most
$\lfloor \log_2 N \rfloor + 1$, so its relative density is zero by \Cref{thm:G13-density} with
$t(N) = N$. The represented set has density one and still misses infinitely many targets.
\end{example}

This is a small arithmetic example of the counting phenomenon, not a thin-orbit result. The setting
that motivates G13 is a group $\Lambda$ acting on a discrete set, with an orbit projected to integers
and observed against a target $A$. There $\Lambda$ is \emph{thin} (Zariski dense of infinite index in
an arithmetic group) and $A$ is \emph{thick} (a finite union of residue classes, so of positive
density).

\begin{schema}[Thin orbits saturate thick targets]\label{thm:G13}
Assume that $\Lambda$ is thin, that its congruence quotients have a uniform spectral gap, that $A$
passes the local congruence conditions, and
\begin{itemize}
  \item \textup{(H1)} an imported theorem gives $m(N) = O(N^{1-\eta})$ for some $\eta > 0$.
\end{itemize}
Then the proportion of missed admissible targets tends to zero.
\end{schema}

\begin{proof}
(H1) gives $m(N) = o(N)$, and a thick target gives $t(N) \ge cN$ for all large $N$, which suffices
for the argument of \Cref{thm:G13-density}. The first three hypotheses are the setting in which such
bounds are proved and are not used.
\end{proof}

Apollonian circle packings are the motivating instance. Curvatures of four mutually tangent circles
satisfy Descartes' relation $k_1^2 + k_2^2 + k_3^2 + k_4^2 = \tfrac12 (k_1 + k_2 + k_3 + k_4)^2$,
and the integral packing generated from $(-1, 2, 2, 3)$ has curvatures only in the residue classes
$\{2, 3, 6, 11, 14, 15, 18, 23\}$ modulo $24$. Bourgain and Fuchs proved that the curvatures have
positive density \citep{BourgainFuchs2010}, and Bourgain and Kontorovich that almost every admissible
integer occurs, with $O(N^{1-\eta})$ exceptions \citep{BourgainKontorovich2014Apollonian}. Haag,
Kertzer, Rickards and Stange showed that the full local-to-global conjecture is nevertheless false
for many packings: certain quadratic and quartic families pass the congruence conditions and are
still missed, for reasons of quadratic reciprocity \citep{HaagEtAl2024}. Those families have density
zero, so they do not contradict the template; they are a second kind of exception, distinct from
failing a congruence condition, and \Cref{ex:G13-powers} is the same pattern in miniature.

\section{General theorems}\label{sec:theorems}

The two entries in this section are theorems whose hypotheses are ordinary data, checked system by
system. Their classical instances rest on imported results.

\subsection{G2: \lawGTwo}\label{sec:G2}

\emph{Name.} The \emph{escrow} is a value, held in a well-founded order, that must fall at every
step; its fall pays for termination. The presentation may change at every step, and the escrow is
compared across the change.

Let $x_0, x_1, \ldots$ be a run with $x_{k+1} = T_k(x_k)$ and terminal set $A$. At each step a
\emph{presentation} $\pi_k : X \to P_k$ records how the state is written, and an \emph{escrow}
$e_k : P_k \to W$ assigns it a value in a well-founded order $(W, <)$.

\begin{theorem}[Transfinite escrow]\label{thm:G2}
Assume $(W, <)$ is well-founded and
\begin{itemize}
  \item \textup{(H1)} whenever $x_k \notin A$,
    $e_{k+1}\bigl(\pi_{k+1}(T_k(x_k))\bigr) < e_k\bigl(\pi_k(x_k)\bigr)$.
\end{itemize}
Then some $x_k$ lies in $A$.
\end{theorem}

\begin{proof}
Otherwise the escrow values $e_k(\pi_k(x_k))$ would form an infinite strictly decreasing sequence in $W$.
\end{proof}

This is Floyd's ranking principle \citep{Floyd1967}. The condition compares the escrow before a
step with the escrow after the step \emph{and} the change of presentation, checked as one
transition. A decrease that holds only before or only after an unrecorded change of presentation
does not discharge it.

Goodstein sequences are the classical instance. In base $k$ the state $n$ is written in hereditary
base $k$; one step rewrites it in base $k+1$ and subtracts $1$. Replacing every base symbol by
$\omega$ gives an ordinal escrow below $\varepsilon_0$ that is unchanged by the rewriting and
strictly decreased by the subtraction. From $4 = 2^2$, rewriting gives $3^3 = 27$ and subtracting
gives $26 = 2\cdot 3^2 + 2 \cdot 3 + 2$: the number grows while the escrow drops from
$\omega^\omega$ to $\omega^2 \cdot 2 + \omega \cdot 2 + 2$. Every Goodstein sequence therefore
terminates \citep{Goodstein1944}, and by the theorem of \citet{KirbyParis1982} this cannot be proved
in Peano arithmetic.

\subsection{G7: \lawGSeven}\label{sec:G7}

\emph{Name.} The angel's \emph{mobility} is set against the devil's power to \emph{confine} it.

In the angel game on $\mathbb{Z}^2$ with the Chebyshev distance, an angel of power $p$ moves each turn
to an undeleted square within distance $p$, and the devil then deletes at most $b$ squares. The angel
wins if it can move forever. The theorem below concerns any game of this shape. A state consists of
the angel's position and a visible obstruction record. At each turn the angel moves to a position
that a legality predicate allows, and the devil then replaces the record by one that a second
legality predicate allows. A strategy is a function of the current state. The angel is
\emph{stuck} at a state with no legal move.

\begin{theorem}[Escape and confinement certificates]\label{thm:G7}
Fix a game of this shape.
\begin{enumerate}
  \item Let $\sigma_A$ be an angel strategy and $I$ a set of states
    such that at every state in $I$ the move of $\sigma_A$ is legal and every legal devil reply to it
    leads to a state in $I$. Let $\sigma_D$ be a devil strategy whose reply to every legal move is
    legal. Then the play of $\sigma_A$ against $\sigma_D$ from any state in $I$ stays in $I$, and
    the angel is never stuck.
  \item Let $\sigma_D$ be a devil strategy and $r$ a
    natural-number rank on states such that the reply of $\sigma_D$ to every legal angel move is
    legal and strictly lowers $r$. Then against every angel strategy that makes a legal move
    whenever one exists, the angel is stuck at some turn $n \le r(s_0)$, where $s_0$ is the
    initial state.
  \item No game has an angel strategy $\sigma_A$ with a set $I$ as in part~1, a devil strategy
    $\sigma_D$ with a rank $r$ as in part~2, and a state in $I$.
\end{enumerate}
\end{theorem}

\begin{proof}
Parts 1 and 2 are inductions on turns, the second bounded by the rank. For part 3, play $\sigma_A$
against $\sigma_D$ from a state in $I$: by part~1 the angel moves legally at every turn, and each
reply lowers the rank, which cannot happen forever.
\end{proof}

The theorem checks certificates; it does not produce strategies. Without a state in $I$ part~3
fails, since an empty set satisfies the hypothesis of part~1 vacuously. For the standard game with
$b = 1$ the answer is known: the angel of power $1$ loses \citep{BerlekampConwayGuy1982,Conway1996},
and an angel of power $2$ wins \citep{Mathe2007,Kloster2007}, with further proofs for power $4$ and
high power \citep{Bowditch2007,Gacs2007}.


\section{Templates}\label{sec:templates}

Each of the four entries in this section is a template: an implication whose hard hypothesis is not
established, in this paper, for any system of interest. The open obligation is named in each case,
and each entry is preceded by what is proved about its principal system. Hypotheses are labelled
(H1), (H2), \dots\ within each entry.

\subsection{G1: \lawGOne}\label{sec:G1}

\emph{Name.} A point is \emph{solvent} once it has earned a certificate of descent; the debt it
carries before that is \emph{hidden} in the gap between the orbit and its start, and paid off
\emph{amortized} over many steps.

\subsubsection*{The concrete system: the Collatz ledger}\label{sec:collatz}

\begin{definition}[Accelerated Collatz map]\label{def:collatz}
For a natural number $n$ let $a(n)$ be the exponent of the largest power of $2$ dividing $3n+1$, and
let $T(n) = (3n+1)/2^{a(n)}$. For a starting value $n$ write $n_j = T^j(n)$ and $a_j = a(n_j)$, and
define $A_k = \sum_{j<k} a_j$ together with $B_0 = 0$ and $B_{j+1} = 3B_j + 2^{A_j}$.
\end{definition}

On odd positive integers $T$ is the usual accelerated map: $T(n)$ is again odd and positive, and
$T(1) = 1$. The definition makes sense for every natural number, and the statements below hold
there; the odd positive integers are the case of interest.

\begin{theorem}[Affine ledger]\label{thm:affine-ledger}
For every natural number $n$ and every $k \ge 0$,
\[
  2^{A_k}\, n_k \;=\; 3^k\, n + B_k .
\]
\end{theorem}

\begin{proof}
For $k = 0$ both sides equal $n$. By definition $2^{a_k} n_{k+1} = 3 n_k + 1$, so
\[
  2^{A_{k+1}} n_{k+1} \;=\; 2^{A_k}\,(3n_k + 1) \;=\; 3\,(3^k n + B_k) + 2^{A_k}
  \;=\; 3^{k+1} n + B_{k+1}. \qedhere
\]
\end{proof}
The ledger is standard in the Collatz literature \citep{Lagarias1985}.

For $n > 0$ the ledger turns descent into an inequality between integers: $n_k < n$ exactly when
$2^{A_k} n > 3^k n + B_k$. This is the certificate that G1 asks for.

\begin{theorem}[Split pairs]\label{thm:split-pair}
Let $m \ge 1$, let $c \in \{1, 2\}$ satisfy $2^{m+1} c \equiv 1 \pmod 3$, and put
\[
  n_1 \;=\; \frac{2^{m+1} c - 1}{3}, \qquad n_2 \;=\; n_1 + 2^m .
\]
Then $n_1 \equiv n_2 \pmod{2^m}$, while $a(n_1) \ge m + 1$ and $a(n_2) = m$.
\end{theorem}

\begin{proof}
The congruence is immediate. By construction $3 n_1 + 1 = 2^{m+1} c$, so $2^{m+1}$ divides
$3n_1 + 1$. Further $3 n_2 + 1 = 2^{m+1} c + 3 \cdot 2^m = 2^m (2c + 3)$, and $2c + 3$ is odd.
\end{proof}

Both numbers are odd and positive, since $3n_1 + 1$ is even and $n_1 \ge 1$. For $m = 1$ the pair
is $(1, 3)$; for $m = 2$ it is $(5, 9)$, the pair of \Cref{sec:intro-example}.

\begin{corollary}\label{cor:no-residue}
For no $m \ge 1$ is the division exponent $a(n)$ of an odd positive $n$ a function of $n$ modulo
$2^m$. In the language of \Cref{sec:setting}, no quotient $n \mapsto n \bmod 2^m$ is
future-sufficient for the ledger.
\end{corollary}

\begin{proof}
The pair of \Cref{thm:split-pair} consists of odd positive numbers, as noted after it; they agree
modulo $2^m$ and have different division exponents.
\end{proof}

\begin{remark}[Three further facts]\label{rem:collatz-paper}
\begin{enumerate}
  \item In the $2$-adic integers, $-1$ is a fixed point: $3(-1) + 1 = -2$ has exponent $1$. It is
    not an integer orbit, but integer orbits can imitate it for a long time.
  \item For $n = 2^L - 1$ with $L \ge 2$, induction gives $n_i = 3^i\, 2^{L-i} - 1$ and $a_i = 1$
    for $0 \le i \le L-2$. The orbit therefore stays above its start for at least $L - 1$ steps,
    so no bound on the first descent time that is uniform in $n$ exists.
  \item If an orbit returns to its start after $k$ steps, then $(2^{A_k} - 3^k)\, n = B_k$. A
    cycle of positive integers can therefore occur only for exponent sequences that make
    $B_k/(2^{A_k} - 3^k)$ a positive integer.
\end{enumerate}
\end{remark}


\subsubsection*{The template}

Let $T : X \to X$ be an update with a set $A$ of terminal states, and let $<$ be a declared
relation on $X$, usually an order. A point $x$ \emph{descends at step $k$} when $T^k(x) < x$. For
$k \ge 1$ the \emph{bad-tail set} $N_k$ consists of the points $x \notin A$ that descend at no step
$j$ with $1 \le j \le k$. The system is \emph{live} when every $x \notin A$ descends at some step.

\begin{proposition}[Bad-tail reformulation]\label{prop:G1}
The system is live exactly when no point lies in every $N_k$:
\[
  \bigl(\forall x \notin A\ \exists k \ge 1 :\ T^k(x) < x\bigr)
  \iff \bigcap_{k \ge 1} N_k = \varnothing .
\]
\end{proposition}

\begin{proof}
A point in every $N_k$ never descends, and a point that never descends lies in every $N_k$.
\end{proof}

This only restates liveness. The content of G1 is in what could empty the intersection. A point in
every $N_k$ has an infinite orbit that never descends. Suppose the system maps into a larger space
$\widehat X$ --- for Collatz, the $2$-adic integers --- whose non-descending fixed or periodic
points outside $X$ form a \emph{ghost set} $\Gamma$, and let $U_1, U_2, \ldots$ be subsets of
$\widehat X$, typically shrinking neighbourhoods of $\Gamma$. The template below uses only the sets
$U_h$, not $\Gamma$ or any topology.

\begin{schema}[Ghost discharge]\label{thm:G1b}
Assume
\begin{itemize}
  \item \textup{(H1)} every orbit that never descends from a point outside $A$
    eventually enters each set $U_h$ and stays there, and
  \item \textup{(H2)} every $x \notin A$ has an $h(x)$ and a $K(x)$ with
    $T^t(x) \notin U_{h(x)}$ for all $t \ge K(x)$.
\end{itemize}
Then $\bigcap_k N_k = \varnothing$, so the system is live.
\end{schema}

\begin{proof}
If $x$ lay in every $N_k$, its orbit would eventually be inside $U_{h(x)}$ by the first hypothesis
and eventually outside it by the second.
\end{proof}

\paragraph{Open obligation.} Both hypotheses, for a system of interest. For the Collatz map
$\Gamma$ contains the $2$-adic fixed point $-1$ (\Cref{rem:collatz-paper}), and both are unproved;
for the template to say something about the integers, separation must be established by arguments
about the integers alone. What is established is the concrete structure of \Cref{sec:collatz}:
descent at step $k$ is the integer inequality $2^{A_k} n > 3^k n + B_k$, and no fixed number of
low-order binary digits predicts the exponents that enter it (\Cref{cor:no-residue}). Orbits of
$2^L - 1$ imitate the ghost for $L - 1$ steps.

\subsection{G6: \lawGSix}\label{sec:G6}

\emph{Name.} A \emph{needle} is an obstruction point; here the run's own history
\emph{generates} them.

Let a run on a countable set $X$ be $x_{n+1} = T(x_n, S_n)$, where $S_n = G(x_0, \ldots, x_n)$ is an
\emph{obstruction set computed from the run's own history}. A point is a \emph{hole} of a finite
window $W \subseteq X$ at time $n$ if it lies in $W$ and has not been visited by time $n$. The
intended law says that if the run is driven back towards the holes of a window faster than its
history adds obstructions there, every hole is eventually visited; and if the obstructions around a
hole grow faster, the hole is permanent. Only the final step of each direction is proved.

\begin{proposition}[Covering and hole formation]\label{thm:G6}
\begin{enumerate}
  \item Let $h(n)$ be the number of holes of a window at time $n$.
    If for every $n$ with $h(n) > 0$ there is an $m > n$ with $h(m) < h(n)$, then for every $n$
    there is an $N \ge n$ with $h(N) = 0$.
  \item If $y$ is unvisited at time $N$, and for every
    $n \ge N$ non-visitation of $y$ at time $n$ implies non-visitation at time $n+1$, then $y$ is
    unvisited at every time $n \ge N$.
\end{enumerate}
\end{proposition}

\begin{proof}
Part~1 is strong induction on $h(n)$; part~2 is induction on $n - N$.
\end{proof}

The hypothesis of part~2 is essentially its conclusion.

\paragraph{Open obligation.} A quantitative comparison, for a system of interest, between how often
the run moves into holes and how fast its history adds obstructions, strong enough to imply the
premises of \Cref{thm:G6}. Two designed variants of Recam\'an's rule satisfy those premises
directly: with jumps $2n$ instead of $n$ only even values are visited, so every odd number is a
permanent hole; with jumps of constant size $1$ the run visits exactly $\{0, \ldots, n\}$ by time
$n$. For Recam\'an's sequence itself --- $a_n = a_{n-1} - n$ when that is positive and unused, and
$a_n = a_{n-1} + n$ otherwise --- whether every positive integer appears is open; $852655$ remains
the smallest missing value in long computations \citep{OEIS-A005132,OEIS-A057167}. It is a
candidate, not an instance. Here \emph{endogenous} means that the run itself produces the
obstruction.

\subsection{G9: \lawGNine}\label{sec:G9}

\emph{Name.} Local defects \emph{evacuate} when they stop accumulating and turn into moving
patterns, that is, into \emph{transport}.

Let a cellular automaton or other local update act on configurations over a lattice with
translations, and fix a reference background. A \emph{defect} is a finite witness that the
configuration differs from the background on a local patch. The \emph{defect census} is a ledger
that follows every defect from time $n$ to time $n+1$. Each defect is continued, annihilated, bound
into a finite composite, absorbed into a transporter, or marked unresolved. Every new defect is
charged to a creation event. A \emph{transporter} is a finite pattern with support $W$, period
$\tau > 0$ and displacement $d$. It is certified by the recurrence $c_{t_0 + (k+1)\tau}(x + d) =
c_{t_0 + k\tau}(x)$ for $x$ in $W + kd$; that equality, not visual recognition, certifies motion.
The transporter is not declared at the start; it appears during the run.

\begin{theorem}[Census bound]\label{thm:G9-census}
Let $u_n$ be the number of unresolved defects at time $n$. If $u_{n+1} \le u_n$ for all $n \ge N$,
then $u_n \le u_N$ for all $n \ge N$. In particular $(u_n)$ is bounded.
\end{theorem}

\begin{proof}
Induction on $n \ge N$.
\end{proof}

\begin{schema}[Evacuation into transport]\label{thm:G9}
Assume that the run carries an exact census whose unresolved count eventually stops growing, and
\begin{itemize}
  \item \textup{(H1)} every defect that stays unresolved long enough is annihilated, bound
    into a stationary composite, or assigned to one of finitely many transporter records, each
    satisfying its recurrence.
\end{itemize}
If some persistent defect is assigned to a transporter with $d \ne 0$, then from some time on every
persistent defect assigned to a transporter is covered by finitely many certified records.
\end{schema}

\begin{proof}
After the absorption time, the records to which persistent defects are assigned form the required
finite family.
\end{proof}

\paragraph{Open obligation.} \textup{(H1)}, which carries all the content; the conclusion
is then a restatement. For the traffic automaton rule~184 the obligation is partly met. On a ring of
$n \ge 2$ cells a \emph{jam} is an occupied cell whose right neighbour is occupied.

\begin{proposition}[Rule 184, partial results]\label{prop:rule184}
On a ring of $n \ge 2$ cells under rule~184:
\begin{enumerate}
  \item a jam-free configuration stays jam-free, and every car advances one cell per step, so from
    any jam-free time on the run satisfies the recurrence with period $1$ and displacement $+1$;
  \item the four-cell configuration $1100$ has one jam, becomes $1010$ after one step, and then
    alternates between $1010$ and $0101$;
  \item the six-cell configuration $001011$ has density $\tfrac12$ and one jam, and after one step it
    still has one jam;
  \item the number of cars is conserved;
  \item for $2 \le n \le 8$, every configuration with at most $n/2$ cars is jam-free after $n$ steps,
    and on four cells after one step.
\end{enumerate}
\end{proposition}

\begin{proof}
Parts (1) to (4) follow from the update rule, a car moving exactly when the cell ahead is empty; part (5) is a
finite check.
(1) In a jam-free configuration every car has an empty cell ahead and moves into it, and no two cars
move into the same cell. (2) and (3) are direct computations. (4) A car leaves a cell only for an
empty right neighbour, and a car enters an empty cell only from its left neighbour, so each move
removes one car from one cell and adds it to the next. (5) is a finite check over all configurations
of each ring.
\end{proof}

\Cref{fig:rule184} shows the runs of parts~2 and~3.

\begin{figure}[htbp]
\centering
\begin{tikzpicture}[x=0.55cm,y=0.55cm,font=\footnotesize]
  \node at (2,1.35) {four-cell ring};
  \foreach \x in {0,1,2,3} \node at (\x+0.5,0.35) {$\x$};
  \foreach \t in {0,1,2,3} \node[left] at (-0.18,-\t-0.5) {$t=\t$};
  \foreach \x/\t in {0/0,1/0,0/1,2/1,1/2,3/2,0/3,2/3}
    \fill[black!80] (\x,-\t) rectangle ++(1,-1);
  \draw[gray] (0,0) grid (4,-4);
  \begin{scope}[xshift=4.8cm]
    \node at (3,1.35) {six-cell ring};
    \foreach \x in {0,1,2,3,4,5} \node at (\x+0.5,0.35) {$\x$};
    \foreach \t in {0,1,2,3} \node[left] at (-0.18,-\t-0.5) {$t=\t$};
    \foreach \x/\t in {2/0,4/0,5/0,0/1,3/1,4/1,1/2,3/2,5/2,0/3,2/3,4/3}
      \fill[black!80] (\x,-\t) rectangle ++(1,-1);
    \draw[gray] (0,0) grid (6,-4);
  \end{scope}
\end{tikzpicture}
\caption{Rule-184 updates on two rings, with time downward, cell indices increasing to the right, and occupied cells filled. The rightmost cell is adjacent to cell $0$ in each ring.}
\label{fig:rule184}
\end{figure}


So low-density rings of up to eight cells evacuate their jams within $n$ steps, while the six-cell
example shows that the jam count need not fall at every update; evacuation on every ring remains
unproved, and so does the assignment of each persistent defect required by \textup{(H1)}. The classical
picture is imported: below density $\tfrac12$ blocked cars evacuate and particles move freely,
above it the holes do, and at density $\tfrac12$ the configuration settles into the alternating
pattern \citep{Fuks1997,Nagel1996,ChowdhurySantenSchadschneider2000}. Langton's ant, whose highway
conjecture is open, is not an instance of any claim here; only its unboundedness is proved
\citep{BunimovichTroubetzkoy1992}.

\subsection{G11: \lawGEleven}\label{sec:G11}

\emph{Name.} Local rules force the absence of every translational symmetry, a \emph{global
anti-symmetry}.

Let $C$ be a finite set of local rules for configurations $x : \mathbb{Z}^d \to A$ over a finite
alphabet $A$, for instance a set of Wang tiles. A nonzero $p \in \mathbb{Z}^d$ is a \emph{period}
of $x$ when $x(v + p) = x(v)$ for all $v$, and $x$ is \emph{aperiodic} when it has no nonzero period.
A \emph{hierarchy} for $C$ consists of scales $L_k \to \infty$ and block types at each scale such
that every admissible configuration decomposes into blocks at every scale, larger blocks being
composed of smaller ones.

\begin{schema}[Hierarchy forces aperiodicity]\label{thm:G11}
Assume
\begin{itemize}
  \item \textup{(H1)} for every admissible configuration $x$ and every nonzero $p$, some
    scale $k$ has block or marker structure in $x$, forced by the rules of $C$, that is not
    invariant under translation by $p$, and
  \item \textup{(H2)} whenever forced structure of $x$ at some scale is not invariant
    under translation by $p$, the vector $p$ is not a period of $x$.
\end{itemize}
Then every admissible configuration is aperiodic.
\end{schema}

\begin{proof}
If an admissible $x$ had a nonzero period $p$, the first hypothesis would give a scale whose forced
structure breaks $p$, and the second would say that $p$ is not a period.
\end{proof}

The mechanism is Robinson's \citep{Robinson1971}. Existence of an admissible configuration is a
separate input: the template is stated for every admissible configuration and is vacuous if there
are none.

\paragraph{Open obligation.} A hierarchy for a given tile set, together with a finite region $R_p$
on which the break of each period is visible. Deciding whether a finite Wang tile set admits any
tiling of the plane is undecidable \citep{Berger1966}, so no algorithm sorts an arbitrary tile set
into periodic, aperiodic or empty. The classical constructions establish the hierarchy for
particular tile sets: Berger's set, Robinson's six tiles \citep{Robinson1971}, Penrose's kites and
darts \citep{Penrose1974}, and the hat and spectre monotiles
\citep{SmithEtAl2023Hat,SmithEtAl2023Spectre}.

\section{What the laws do not give}\label{sec:scope}

Each law is easy to read more strongly than it is stated. \Cref{tab:scope} lists the inferences we
expect a reader to be tempted by, and what blocks each one. A complete register of non-implications,
with a separating example for each, is in \Cref{supp:nonimpl}.

{\small
\begin{longtable}{@{}>{\raggedright\arraybackslash}p{5.6cm}>{\raggedright\arraybackslash}p{8.6cm}@{}}
\caption{Tempting inferences and what blocks them.}\label{tab:scope}\\
\toprule
\textbf{Tempting inference} & \textbf{Why it fails} \\
\midrule
\endfirsthead
\toprule
\textbf{Tempting inference} & \textbf{Why it fails} \\
\midrule
\endhead
\bottomrule
\endfoot
The ghost-discharge template says something about the Collatz conjecture. & Both of its hypotheses
are unproved for the Collatz map (\Cref{sec:G1}). The split pairs of \Cref{thm:split-pair} rule out
fixed-depth prediction; they are not evidence about liveness. \\
A decreasing escrow may decrease before or after a change of presentation. & The comparison must
cover the step and the change of presentation together (\Cref{sec:G2}). \\
A statistic that tracks cost is a potential. & Only a nonnegative potential entering the
telescoping identity bounds cost (\Cref{thm:G3}). \\
Covering every target so far shows the cover is exhaustive. & Exhaustion is the hypothesis of
\Cref{thm:G4}; for Erd\H{o}s--Straus it is the open problem. \\
A closed target-free family for base $2$ suggests one for $196$ in base $10$. & None is known, and
the base-$2$ family was found by hand for one orbit (\Cref{sec:G5}). \\
A long run without reaching $852655$ shows it is a permanent hole of Recam\'an's sequence. & A
finite computation says nothing about later steps; Recam\'an's sequence is a candidate, not an
instance (\Cref{sec:G6}). \\
Return pressure exceeding obstruction growth forces coverage. & The implication from the rate
conditions to the premises of \Cref{thm:G6} is open (\Cref{sec:G6}). \\
The certificate theorem for the angel game finds winning strategies. & It checks given strategies
(\Cref{thm:G7}); the thresholds are imported. \\
Confluent move systems have route-independent counters. & \Cref{ex:G8-b} is confluent and its
counter vectors differ. \\
Any terminating move system has a unique final configuration. & \Cref{ex:G8-d} terminates without
one. \\
Visible motion of a pattern is transport. & Only a checked recurrence with displacement certifies
it (\Cref{sec:G9}); Langton's ant is not an instance. \\
Persistence for all iterates can be observed. & It must be proved from protection and regeneration
(\Cref{thm:G10}); a finite run gives bounded evidence only. \\
Local rules whose finite patches extend always force aperiodicity. & A hierarchy must be supplied
(\Cref{thm:G11}), and whether a tile set tiles at all is undecidable. \\
Every impossible global colouring has a finite witness. & Only under the compactness hypothesis
(\Cref{prop:G12-compact}), which uses the axiom of choice. \\
Passing every congruence condition means an integer occurs in the orbit. & Reciprocity
obstructions are missed despite passing (\Cref{sec:G13}). \\
\end{longtable}
}

\section{Formalization}\label{sec:lean}

\paragraph{What is checked.} The general statements of all thirteen entries are checked in Lean~4
\citep{DeMouraUllrich2021}, in a library with no external dependencies, together with the Collatz
ledger and split pairs (\Cref{thm:affine-ledger,thm:split-pair}), the four-phase family and the
orbit of $22$ (\Cref{thm:four-phase,cor:orbit-22}), Ducci nilpotency (\Cref{thm:ducci}), the two
separating examples for G8, and the instances in \Cref{ex:G4-intervals,ex:G10-checksum,ex:G13-powers}
and \Cref{prop:rule184}, the binary counter of \Cref{ex:G3-counter}, a three-site sandpile, the
colouring obstruction of the Moser spindle and the integer Ducci bound for $n = 4$. The last column of
\Cref{tab:catalog} says, entry by entry, what is covered; sandpiles on arbitrary graphs, the
plane embedding of the Moser spindle, the other remarks and the imported results are proved on paper
only. Some Lean statements are formally weaker than the printed ones, for
example \Cref{thm:G2} without the extraction of a terminal stage, and some carry a hypothesis they
do not use; both are listed in \Cref{supp:status}.

\paragraph{Axioms.} The library builds without \texttt{sorry} or added axioms and uses at most the
standard axioms \texttt{propext}, \texttt{Quot.sound} and \texttt{Classical.choice}.

\paragraph{Where to look.} Following \citet{Pollack1998}, what a reader must trust is the kernel and
the statements. \Cref{supp:concordance} maps every numbered result of this paper to its Lean
declarations, file and line; \Cref{supp:status} gives the axioms of each declaration and a
clause-by-clause comparison with the printed statements; \Cref{supp:replay} gives the commands that
rebuild the library.


\clearpage
\bibliographystyle{plainnat}
\bibliography{references}

\end{document}
